\documentclass[10pt,a4paper]{amsart}
\usepackage[margin=19mm]{geometry}
\usepackage{amsmath,amssymb,mathtools}
\usepackage[scr=rsfs,cal=euler]{mathalfa}
\usepackage{xfrac}
\usepackage[unicode,hidelinks,bookmarks=false]{hyperref}
\newcommand{\ie}{i.e.\ }
\newcommand{\cf}{cf.\ }
\newcommand{\resp}{respectively\ }
\newcommand{\Subsection}{\subsection}
\providecommand{\nobreakdash}{\nobreak\hbox{-}}
\setlength{\emergencystretch}{2em}
\multlinegap0pt
\usepackage{amssymb}
\newtheorem{lemma}{Lemma}
\newcommand{\CAT}{\textup{CAT}}
\def\R{\mathbb R}
\providecommand{\abs}[1]{\left\lvert#1\right\rvert}
\newcommand{\ip}[2]{\left\langle #1, #2 \right\rangle}
\title{The proximal point method and its two variants for monotone vector fields in Hadamard spaces}
\author{Parin Chaipunya, Fumiaki Kohsaka}
\date{Source: arXiv:2605.01354v1 (CC BY 4.0)}
\begin{document}
\maketitle

\noindent\emph{Source and licence.} The proximal point method and its two variants for monotone vector fields in Hadamard spaces. Version arXiv:2605.01354v1 (CC BY 4.0), distributed by arXiv under the Creative Commons Attribution 4.0 International licence, http://creativecommons.org/licenses/by/4.0/, as recorded in the arXiv licence metadata for this exact version and shown on the abstract page https://arxiv.org/abs/2605.01354v1. Full licence notice: \texttt{licenses/arxiv-2605.01354-CC-BY-4.0.txt}.\par\smallskip

\noindent\textit{Editorial scope.} Counted unit: the article's lemma lem:d\_p-pm (source0.tex lines 868--873). The article's complete original proof of it is delivered with it (source0.tex lines 875--1044, one \texttt{proof} environment of 170 lines). No mathematics is written, repaired or re-derived: the statement and the proof are the article's own lines, in the article's own notation, and no retained line is substituted or re-flowed.  Retained uncounted as the source-local context the proof needs: lines 235--251; lines 829--831; lines 850--855; lines 858--862.  Nothing was omitted from any retained span, so the package carries no omission record.  No retained line contains a citation, so the wrapper carries no bibliography and no bibliography entry.  No retained line contains a figure environment, an image-inclusion command or drawing code, so no image is delivered and the package declares no asset dependency.  Editorial formatting only, in addition: an AMS \texttt{amsart} preamble that restates the article's own declarations and macros used by the retained lines, namely the environments \texttt{lemma} on separate counters, together with the standard packages those lines need.  The retained environments print in this document as Lemma 1 in order of appearance, whereas the article prints the same items with the numbering of its own full document.  The complete native source of the article as distributed in the arXiv e-print for this exact version is bundled unchanged as \texttt{sources/199632/source0.tex}.
\begin{lemma}\label{lem:Alex}
 Let $x,y,z$ be nonzero points in $\R^2$ 
% Let $x,y,z$ be distinct nonzero points in $\R^2$, 
% and $u$ an element in $\R^2$ which is orthogonal to $y$. 
% Suppose that $\ip{x}{u}\ip{z}{u}\leq 0$ 
 and let 
 \[
  \theta = \arccos \frac{\ip{x}{y}}{\abs{x}\abs{y}} 
 \quad \textrm{and} \quad 
  \theta' = \arccos \frac{\ip{z}{y}}{\abs{z}\abs{y}}. 
 \]
 If $\theta+\theta' \geq \pi$, 
 then 
 \[
  \abs{x-y} + \abs{y-z} \geq \abs{x} + \abs{z}. 
 \]
\end{lemma}

\begin{align}\label{eq:def-Y}
 Y=[0,\infty)\times \Gamma_p, 
\end{align}

\begin{align}\label{eq:def-d_p-L}
 \tilde{d}_p\bigl((\alpha, \gamma), (\beta, \eta)\bigr) 
 =\sqrt{\bigl(\alpha L_{\gamma}\bigr)^2 
 +\bigl(\beta L_{\eta}\bigr)^2 
 -2\bigl(\alpha L_{\gamma}\bigr)\bigl(\beta L_{\eta}\bigr) \cos \angle_p(\gamma, \eta)}  
\end{align}

\begin{align}\label{eq:ineq-d_p}
 \abs{\alpha L_{\gamma} - \beta L_{\eta}} 
 \leq \tilde{d}_p\bigl((\alpha, \gamma), (\beta, \eta)\bigr) 
 \leq \alpha L_{\gamma} + \beta L_{\eta}
\end{align}

\begin{lemma}\label{lem:d_p-pm}
 Let $(X, \rho)$ be a $\CAT(0)$ space, 
 $p$ an element in $X$, and $Y$ the set defined by~\eqref{eq:def-Y}. 
 Then $\tilde{d}_p$ is a pseudometric 
 on $Y$. 
\end{lemma}

\begin{proof}
 It follows from~\eqref{eq:ineq-d_p} that 
 \[
  \tilde{d}_p\bigl((\alpha, \gamma), (\beta, \eta)\bigr)\in [0,\infty)
 \]
 for all $(\alpha, \gamma), (\beta, \eta)\in Y$. 
 Since $\angle_p$ is a pseudometric on $\Gamma_p$, 
 we know that 
 \[
  \tilde{d}_p\bigl((\alpha, \gamma), (\alpha, \gamma)\bigr) =0
 \]
 for all $(\alpha, \gamma)\in Y$ and 
 $\tilde{d}_p$ is symmetric. 
 
 We next show that $\tilde{d}_p$ satisfies the triangle 
 inequality. Let $(\alpha_i, \gamma_i) \in Y$ 
 for $i=1,2,3$. 
 If $\alpha_1 L_{\gamma_1}=0$, 
 then it follows from~\eqref{eq:def-d_p-L} and~\eqref{eq:ineq-d_p} that 
 \begin{align*}
  \tilde{d}_p\bigl((\alpha_1, \gamma_1), (\alpha_3, \gamma_3)\bigr) 
  &=\alpha_3 L_{\gamma_3} \\
  &\leq \alpha_2 L_{\gamma_2} + \abs{\alpha_3 L_{\gamma_3} - \alpha_2 L_{\gamma_2}} \\
  &\leq 
   \tilde{d}_p\bigl((\alpha_1, \gamma_1), (\alpha_2, \gamma_2)\bigr) 
  +
  \tilde{d}_p\bigl((\alpha_2, \gamma_2), (\alpha_3, \gamma_3)\bigr). 
 \end{align*}
 If $\alpha_3 L_{\gamma_3}=0$, then 
 \begin{align*}
  \tilde{d}_p\bigl((\alpha_1, \gamma_1), (\alpha_3, \gamma_3)\bigr) 
  &=\alpha_1 L_{\gamma_1} \\
  &\leq \abs{\alpha_1 L_{\gamma_1} - \alpha_2 L_{\gamma_2}} + \alpha_2 L_{\gamma_2} \\
  &\leq 
   \tilde{d}_p\bigl((\alpha_1, \gamma_1), (\alpha_2, \gamma_2)\bigr) 
  +
  \tilde{d}_p\bigl((\alpha_2, \gamma_2), (\alpha_3, \gamma_3)\bigr). 
 \end{align*}
 If $\alpha_2 L_{\gamma_2} =0$, then we have 
 \begin{align*}
  \tilde{d}_p\bigl((\alpha_1, \gamma_1), (\alpha_3, \gamma_3)\bigr) 
  &\leq \alpha_1 L_{\gamma_1} + \alpha_3 L_{\gamma_3} \\
  &=
   \tilde{d}_p\bigl((\alpha_1, \gamma_1), (\alpha_2, \gamma_2)\bigr) 
  +
  \tilde{d}_p\bigl((\alpha_2, \gamma_2), (\alpha_3, \gamma_3)\bigr). 
 \end{align*}
 We next consider the case where $\alpha_i L_{\gamma_i} >0$ 
 for all $i=1,2,3$. Set $\theta=\angle_p(\gamma_1, \gamma_2)$ 
 and $\theta'=\angle_p(\gamma_2, \gamma_3)$. 
 Let $\overline{y}_1, \overline{y}_2, \overline{y}_3$ 
 be points in $\R^2$ given by 
 \[
  \overline{y}_1 = 
   \alpha_1 L_{\gamma_1}
  \begin{pmatrix}
   1 \\
   0
  \end{pmatrix}, 
  \quad 
  \overline{y}_2 = 
  \alpha_2 L_{\gamma_2}
  \begin{pmatrix}
   \cos \theta  \\
   \sin \theta 
  \end{pmatrix}, 
  \quad \textrm{and} \quad 
  \overline{y}_3 = 
   \alpha_3 L_{\gamma_3}
   \begin{pmatrix}
   \cos (\theta + \theta') \\
   \sin (\theta + \theta')
  \end{pmatrix}. 
 \] 
 Then we have
 \begin{itemize}
  \item $\abs{\overline{y}_i} = \alpha_i L_{\gamma_i}$ 
 for all $i=1,2,3$; 
  \item $\arccos \frac{\ip{\overline{y}_1}{\overline{y}_2}}
 {\abs{\overline{y}_1}\abs{\overline{y}_2}} = \theta$ 
 and $\arccos \frac{\ip{\overline{y}_2}{\overline{y}_3}}
 {\abs{\overline{y}_2}\abs{\overline{y}_3}} = \theta'$. 
%  \item $\ip{\overline{y}_1}{\overline{u}}\ip{\overline{y}_2}{\overline{u}}\leq 0$, 
% where $\overline{u}$ denotes a unit vector in $\R^2$ 
% such that $\ip{\overline{y}_2}{\overline{u}}=0$. 
 \end{itemize}
 Consequently, we have 
 \begin{align}
  \begin{split}\label{eq:d_p-pm-a}
  &\tilde{d}_p\bigl((\alpha_1, \gamma_1), (\alpha_2, \gamma_2)\bigr) \\
  &= \sqrt{\bigl(\alpha_1 L_{\gamma_1}\bigr)^2 
 +\bigl(\alpha_2 L_{\gamma_2}\bigr)^2 
 -2\bigl(\alpha_1 L_{\gamma_1}\bigr)
 \bigl(\alpha_2 L_{\gamma_2}\bigr)\cos \angle_p(\gamma_1, \gamma_2)} \\
  &= \sqrt{\abs{\overline{y}_1}^2 + \abs{\overline{y}_2}^2 
 -2\abs{\overline{y}_1}\abs{\overline{y}_2}
 \cos \theta} \\
  &= \sqrt{\abs{\overline{y}_1}^2 + \abs{\overline{y}_2}^2 
 -2\ip{\overline{y}_1}{\overline{y}_2}} 
  =\abs{\overline{y}_1-\overline{y}_2}. 
  \end{split}
 \end{align}
 Similarly, we have 
 \begin{align}
  \begin{split}\label{eq:d_p-pm-b}
  &\tilde{d}_p\bigl((\alpha_2, \gamma_2), (\alpha_3, \gamma_3)\bigr) \\
  &= \sqrt{\bigl(\alpha_2 L_{\gamma_2}\bigr)^2 
    +\bigl(\alpha_3 L_{\gamma_3}\bigr)^2 
 -2\bigl(\alpha_2 L_{\gamma_2}\bigr)
 \bigl(\alpha_3 L_{\gamma_3}\bigr)\cos \angle_p(\gamma_1, \gamma_3)} 
  = \abs{\overline{y}_2-\overline{y}_3}.       
  \end{split}
 \end{align}

 If $\theta+\theta'\leq \pi$, then we have 
 \[
  \theta + \theta' 
  = \arccos \frac{\ip{\overline{y}_1}{\overline{y}_3}}
  {\abs{\overline{y}_1}\abs{\overline{y}_3}}.  
 \]
 By the triangle inequality for $\angle_p$, we have 
 \[
  \angle_p(\gamma_1, \gamma_3) 
  \leq \angle_p(\gamma_1, \gamma_2) +   \angle_p(\gamma_2, \gamma_3) 
  =\theta + \theta' =
  \arccos \frac{\ip{\bar{y}_1}{\bar{y}_3}}{\abs{\bar{y}_1}\abs{\bar{y}_3}}
 \]
 and hence it follows from 
 the triangle inequality 
 for $\abs{\,\cdot \,}$,~\eqref{eq:d_p-pm-a}, 
 and~\eqref{eq:d_p-pm-b} that 
 \begin{align*}
  \begin{split}
  \tilde{d}_p\bigl((\alpha_1, \gamma_1), (\alpha_3, \gamma_3)\bigr) 
  &=\sqrt{\abs{\overline{y}_1}^2 + \abs{\overline{y}_3}^2 
 -2\abs{\overline{y}_1}\abs{\overline{y}_3}\cos \angle_p(\gamma_1, \gamma_3)} \\
  &\leq \sqrt{\abs{\overline{y}_1}^2 + \abs{\overline{y}_3}^2 
 -2\abs{\overline{y}_1}\abs{\overline{y}_3}
 \cos \left(
 \arccos \frac{\ip{\bar{y}_1}{\bar{y}_3}}{\abs{\bar{y}_1}\abs{\bar{y}_3}}
 \right)
 } \\
  &= \sqrt{\abs{\overline{y}_1}^2 + \abs{\overline{y}_3}^2 
 -2\ip{\bar{y}_1}{\bar{y}_3}} \\
  &=\abs{\overline{y}_1-\overline{y}_3} \\
  &\leq \abs{\overline{y}_1-\overline{y}_2} + \abs{\overline{y}_2-\overline{y}_3} \\
  &= \tilde{d}_p\bigl((\alpha_1, \gamma_1), (\alpha_2, \gamma_2)\bigr) 
  + \tilde{d}_p\bigl((\alpha_2, \gamma_2), (\alpha_3, \gamma_3)\bigr). 
  \end{split}
 \end{align*}
On the other hand, if $\theta + \theta' > \pi$, 
 then it follows from Lemma~\ref{lem:Alex} that 
 \begin{align}\label{eq:d_p-pm-c}
  \abs{\overline{y}_1 - \overline{y}_2} + \abs{\overline{y}_2-\overline{y}_3} 
  \geq \abs{\overline{y}_1} + \abs{\overline{y}_3}. 
 \end{align}
 Then it 
 follows from~\eqref{eq:ineq-d_p},~\eqref{eq:d_p-pm-a},~\eqref{eq:d_p-pm-b}, and~\eqref{eq:d_p-pm-c} that 
 \begin{align*}
  \begin{split}
   \tilde{d}_p\bigl((\alpha_1, \gamma_1), (\alpha_2, \gamma_2)\bigr) 
  + \tilde{d}_p\bigl((\alpha_2, \gamma_2), (\alpha_3, \gamma_3)\bigr) 
  &= \abs{\overline{y}_1 - \overline{y}_2} + \abs{\overline{y}_2-\overline{y}_3} \\
  &\geq \abs{\overline{y}_1} + \abs{\overline{y}_3} \\
  &= \alpha_1 L_{\gamma_1} + \alpha_3 L_{\gamma_3} \\
  &\geq \tilde{d}_p\bigl((\alpha_1, \gamma_1), (\alpha_3, \gamma_3)\bigr). 
  \end{split}
 \end{align*}
 Therefore, $\tilde{d}_p$ is a pseudometric on $Y$. 
\end{proof}

\end{document}
